\documentclass[11pt,letterpaper]{article}
\usepackage[utf8]{inputenc}
\usepackage[T1]{fontenc}
\usepackage[spanish,es-nodecimaldot,es-tabla]{babel}
\usepackage{mathptmx}
\usepackage{amsmath,amssymb}
\usepackage{graphicx}
\usepackage{booktabs,array}
\usepackage{tabularx}
\usepackage{float}
\usepackage{enumitem}
\usepackage[hyphens]{url}
\usepackage[hidelinks]{hyperref}
\usepackage{geometry}
\usepackage{xcolor}
\usepackage[most]{tcolorbox}
\usepackage{pdflscape}
\geometry{letterpaper,left=2.4cm,right=2.4cm,top=2.3cm,bottom=2.3cm}
\linespread{1.1}
\setlength{\parindent}{0pt}
\setlength{\parskip}{0.5em}
\setlist[itemize]{leftmargin=1.4em,topsep=1pt,itemsep=1pt}
\setlist[enumerate]{leftmargin=1.7em,topsep=1pt,itemsep=1pt}
\definecolor{iiBlue}{HTML}{1F4E79}
\definecolor{iiOrange}{HTML}{C55A11}
\newtcolorbox{propbox}[1]{colback=iiBlue!6,colframe=iiBlue,coltitle=white,fonttitle=\bfseries,title=#1,arc=2pt,boxrule=0.9pt,left=6pt,right=6pt,top=3pt,bottom=3pt}
\newtcolorbox{alertbox}[1]{colback=iiOrange!7,colframe=iiOrange,coltitle=white,fonttitle=\bfseries,title=#1,arc=2pt,boxrule=0.9pt,left=6pt,right=6pt,top=3pt,bottom=3pt}
\newcommand{\hrulethin}{\noindent\rule{\textwidth}{0.6pt}}
\begin{document}\pagestyle{plain}
\begin{center}
    \textbf{\large Universidad de Santiago de Chile}\\[0.1cm]
    Departamento de Ingeniería Industrial
\end{center}
\vspace{0.15cm}\hrulethin\vspace{0.35cm}
\begin{center}
    {\scshape Estadística Aplicada \quad$\cdot$\quad Caso 2 · Clasificación supervisada}\\[0.35cm]
    {\LARGE \textbf{Guía de presentación}}\\[0.15cm]
    {\large Proyecto Integrador: Clasificación Supervisada}\\[0.2cm]
    Predicción de una respuesta categórica
\end{center}
\vspace{0.2cm}\hrulethin\vspace{0.4cm}
\section*{Condiciones}
Diez diapositivas de contenido, \textbf{10 minutos de exposición} y \textbf{10 minutos de preguntas}. Participan
todos los integrantes; no hay roles fijos y cualquiera puede ser consultado sobre cualquier parte. Las
diapositivas de título y respaldo no cuentan dentro de las diez.

\begin{alertbox}{Lo que más pesa}
La evaluación se concentra en la \textbf{oratoria y el dominio conceptual} y en la \textbf{calidad de la presentación}. Exponer leyendo o no manejar los conceptos baja sustancialmente la nota.
\end{alertbox}

\section*{Contenido esperado (10 láminas)}
\begin{center}\small
\begin{tabular}{c p{13.2cm}}
\toprule
Slide & Contenido esperado \\
\midrule
1 & Qué se predice, quién actúa y por qué los errores difieren. \emph{Problema y usuario:} Contexto y relevancia, Usuario de la predicción, Decisión apoyada, Unidad de observación, Población objetivo. \emph{Clases y costos:} Clases y clase positiva, Instante y horizonte, Significado de FP y FN, Costo relativo, Capacidad. \\
2 & La elección se justifica por la decisión, no por disponibilidad. \emph{Fuente:} Fuente y licencia, Período y tamaño, Población. \emph{Pertinencia:} Variables al corte, Prevalencia, Alternativa descartada. \\
3 & Hallazgos que modificaron preparación, partición o métrica. \emph{Calidad:} Faltantes y duplicados, Dependencia, Variables sospechosas. \emph{Desbalance y fuga:} Prevalencia por grupo, Subgrupo crítico, Riesgos de leakage. \\
4 & Mismos resamples, test intacto, desbalance dentro de folds. \emph{Partición:} Temporal o por grupos, Test intacto. \emph{Resamples y desbalance:} Folds comunes y semillas, Pesos o SMOTE dentro de folds. \\
5 & Qué aporta cada modelo y qué hiperparámetros la controlan. \emph{Baseline y lineales:} Baseline operacional, Logística, KNN. \emph{Margen y árboles:} SVM lineal y kernel, Árbol y Random Forest, Boosting. \\
6 & Contraste baseline, variabilidad y prevalencia. \emph{ROC-AUC / PR-AUC:} Frente al baseline, Variabilidad entre folds. \emph{Métricas de clase:} Precision y recall, F1 y especificidad. \\
7 & Si la política usa probabilidades, evalúe calibración. \emph{Diagnóstico:} Log-loss o Brier, Curva de calibración. \emph{Consecuencia:} Ranking vs probabilidad, Efecto en la decisión. \\
8 & Seleccione umbral o top-$k$ por costos y capacidad. \emph{Criterio:} Costo esperado, Capacidad o top-k, Elegido en validación. \emph{Punto de operación:} Matriz de confusión, Volumen de acciones. \\
9 & Una sola evaluación de test del sistema congelado. \emph{Test:} Sistema congelado, Resultado único. \emph{Subgrupo crítico:} Errores frecuentes, Denominadores. \\
10 & Decisión accionable con límites explícitos. \emph{Recomendación:} Qué usar y quién, Población y límites. \emph{Monitoreo:} Indicadores y drift, Cuándo recalibrar. \\
\bottomrule
\end{tabular}\end{center}

\section*{Criterios de calidad}
\begin{itemize}
\item Una idea central por lámina, expresada como conclusión visible.
\item Exponer, no leer: las láminas apoyan el discurso, no lo reemplazan.
\item Toda figura con título, unidades, leyenda legible y un comentario interpretativo.
\item Sin volcados de código ni tablas con exceso de decimales.
\item Mostrar la comparación frente al baseline o a la referencia, no una sola cifra.
\item Coherencia exacta entre presentación, script e informe.
\end{itemize}
\section*{Qué baja la nota}
\begin{itemize}
\item Leer gran parte de la exposición.
\item No manejar los conceptos al ser preguntado.
\item Presentar resultados sin baseline, sin incertidumbre o con fuga de información.
\item Afirmar causalidad desde importancia de variables o asociaciones.
\end{itemize}
\section*{Defensa}
Cualquier integrante puede ser consultado sobre cualquier parte. Prepare respaldo sobre datos, validación, hiperparámetros y decisiones.
\end{document}
